\section{The finite-horizon information problem}\label{sec:scope}
An exact likelihood and an exact state for a prescribed future decision
need not have the same dimension. We study this distinction for experiments
in which observation probabilities are polynomial functions of a static
unknown parameter. The observation channel, the allowed interventions and
the resource count are specified together. In particular, an unsuccessful
preparation or a rejected report remains an observation and consumes one
trial. This makes the attainable likelihood family a consequence of a
calibrated apparatus, not a collection of parameter-dependent commands.

The basic laboratory has a radius $R$ in $I=[9/20,47/100]$. A cartridge is
one parameter-independent ambient placement attempt, with independent
prepared pointer and comparator coordinates. Its collision window is the
physical time $T$, and its raw polynomial kernel is derived by integrating
all one-collision tubes. The full circular detector permits the compact
mode family
\[
 \cA=[T_-,T_+]\times[0,E]\times\mathbb S^1,
 \qquad 0<T_-\le T_+<3/50,\qquad E(47/100)^2<1.
\]
The comparator may use $0\le g\le1$. The central finite detector below
instead fixes one nonzero-amplitude mode, discards the angle except for a
single collision bit, and permits only four lookup probabilities in
$[\eta,1-\eta]$ with $0<\eta<1/2$. Its accessible alphabet and firmware are
fixed in advance. Internally measured collision coordinates are not
available to the controller. Both instruments use the same preparation;
neither obtains an uncharged equilibrium reset.

\subsection{The state and decision law}
For a positive finite-alphabet polynomial experiment of degree $q$, assume
that the finite lookup comparator has an open family of attainable failure
factors in $\mathcal P_q$. Fix a full-support prior on the compact parameter
interval. After $n$ reports a normalized positive polynomial of degree at
most $qn$ represents the full posterior. Consider instead decisions about
$m$ further reports. Operational equivalence means equality of their laws
for every common continuation policy; the future task is the same, and no
terminal payoff in this definition directly inspects the unknown parameter.

\begin{theorem}[Attainable finite-horizon state and decision law]\label{thm:decision-main}
Under these assumptions, for $n,m\ge1$:
\begin{enumerate}
\item Operational equivalence is exactly equality of the posterior moments
through degree $qm$. The attainable quotient has dimension
$d=q\min(n,m)$, contains a relatively open $d$-dimensional patch, and admits
an exact continuous state with $d$ coordinates. No smaller continuous
Euclidean state is exact. For total budget $N$ these states can be updated
causally with dimension $q\min(n,N-n)$ at time $n$.
\item There is a fixed observable prediction-score experiment with exactly
$m$ further cartridges for which the optimal worst-history regret using
$M$ stored symbols is bounded above and below by positive constants times
$M^{-2/d}$. The same exponent holds for unconditional expected regret in
one fixed independent uniform-command exploration experiment. The comparison includes discontinuous encoders and fresh
randomization, but no uncharged tape of past reports or commands. Equivalently,
$B$ mutable bits give the sharp exponent $2^{-2B/d}$.
\item A fixed four-symbol Lorentz detector realizes $q=3$ with four-entry
lookup gates. Erasing its collision sign before reporting and comparison
realizes $q=2$. Under the same four-entry command convention, query words,
prior, score and cartridge counts, sign erasure changes $d$ from
$3\min(n,m)$ to $2\min(n,m)$ and changes the memory exponent accordingly.
Retaining the sign also yields a strictly positive two-cartridge prediction
gain with zero acceptance cost, given by \eqref{eq:zero-cost-gain}.
\end{enumerate}
\end{theorem}
\begin{proof}
Theorem~\ref{thm:calibrated-rank} calculates the finite calibration matrix
and its exact cube image. Lemma~\ref{lem:probe-span} realizes a spanning
family of polynomial tests as actual all-failure events. Theorem
\ref{thm:decision-quotient} identifies their operational quotient by the
positive polynomial Gram form, and Theorem~\ref{thm:causal-horizon} gives
its causal implementation. Lemma~\ref{lem:score-metric} turns forecast
regret into squared distance on that quotient. Theorem~\ref{thm:memory-rate}
proves matching finite-code bounds; Theorem~\ref{thm:average-memory-rate}
proves their unconditional fixed-experiment counterpart; and
Theorem~\ref{thm:continuous-regret}
gives a quantitative obstruction for smaller continuous real states.
The common-menu construction makes the ablation a comparison of the same
score protocol; Theorem~\ref{thm:zero-cost-gain} proves the strict gain.
\end{proof}

The constants depend on the fixed prior, detector and checkpoint. In
particular, the lower constant is not uniform as the collision amplitude
vanishes. A full-support prior need not have a Lebesgue density. Full
support is used because the polynomial Gram form is then positive definite.
The memory bound concerns the mutable history-dependent state; immutable
calibration data and computation of a codebook are not asserted to have
zero physical or algorithmic cost.

\subsection{Relation to the complete likelihood and control framework}
The $qn$ full-likelihood dimension is retained, not contradicted, by the
smaller observable-future quotient. If a terminal decision has an arbitrary
continuous payoff $G(R,d)$ depending on the hidden radius, the complete
positive likelihood remains the state used by the control theorem. A
specified polynomial parameter payoff admits the additional moment bound
in Section~\ref{sec:decision-horizon}; its nominal degree alone does not
prove a matching lower bound for a restricted decision class.

For the full circular detector we retain the general measurable-comparator
realizability theorem, positive Bernstein filter, compact gate-moment
optimization, boundary-complete attained feedback, quantitative lookup
approximation, and response of every fixed parameter-independent controller.
The collision high-acceptance region can be chosen as at most two arcs at
every stage under tagwise current accepted rewards. A normalized positive
Bernstein surrogate treats smooth nonpolynomial marks with complete-policy
value and regret bounds. These are statements about specified readouts;
they do not identify a singular complete collision record with a positive
detector density. The exact whole-preparation and hybrid response arguments
are retained in the foundation companion, with the current qualifications
in Appendix~\ref{app:repairs}.

The later two-step adaptive cost comparison is also retained with its
complete full-Borel-class reduction and rational certificate. Its strict
advantage does not require the collision bit: the optimal nominal policy
uses only no collision versus its complement. The central result here is
instead a sharp state--memory--decision law whose exponent changes when
that bit is removed, together with an independently positive zero-cost
prediction gain. No statement about a long uninterrupted billiard orbit
is inferred from independent short-window cartridges.

\subsection{Established mechanisms and the mathematical distinction}
Action-conditional tests and predictive state representations are established
\cite{LittmanSutton,SinghJamesRudary}. Our state is a finite-horizon nonlinear
quotient of attainable likelihoods for a static continuous parameter. We
compute its dimension and a corresponding memory--decision law; we do not
claim a horizon-independent linear rank for an infinite system-dynamics
matrix. Bernstein multiplication and beta-mixture representations likewise
use established polynomial arithmetic and statistical constructions
\cite{FaroukiRajan,PetroneWasserman}.

Convex continuation envelopes are familiar in partially observed control
\cite{SmallwoodSondik,FeinbergEtAl}. Squared-error quantization supplies the
covering and volume mechanism behind the exponent $2/d$ \cite{GrafLuschgy},
and Borsuk--Ulam supplies the topological obstruction for a smaller
continuous real state \cite{Matousek}. Comparison of experiments and
conditional variance explain the nonnegative value of an extra observation
\cite{Blackwell}. The distinct assertion is that one explicitly constrained,
counted experiment has a computed attainable quotient, a causal exact
state, a matching observable-score memory law, and a same-protocol physical
ablation that changes that law. The proof includes the apparatus-to-factor
map and both directions of the state and decision bounds.

The fixed-space rank extension and the frozen-failure continuation
argument are credited to the repository referee report \cite{RefereeV4}
where they are used. The general flux and saltation mechanisms are classical
\cite{Golse,Chernov,KongEtAl} and are rederived to control normalization,
failed preparation, and the fixed reporting space. The source record and
execution receipts are separate from the mathematical arguments.
